\chapter{Valutazione Sperimentale}\label{chap:sixth}

\inlineminitoc

Questo capitolo misura le due strategie di recupero definite nel Capitolo~\ref{chap:fifth} e risponde alla domanda che motiva l'intero lavoro: \textit{la struttura esplicita di un grafo normativo migliora il recupero rispetto alla sola similarità semantica, e a quale prezzo?}

La Sezione~\ref{sec:antitesi} verifica infine se il contributo del grafo persista quando la ricerca vettoriale usa rappresentazioni già contestualizzate a livello di articolo.

\section{Il banco di domande}
\label{sec:banco}

Non esiste, per il corpus normativo impiegato in questo lavoro, un insieme di domande e risposte di riferimento. È stato quindi necessario costruirne uno, e la sua composizione è una scelta metodologica che condiziona ogni numero riportato nel seguito.

\subsection{Struttura del banco}

Il benchmark diagnostico è costituito da \textbf{70 domande}, organizzate in \textbf{sette categorie da dieci}. Ogni domanda dichiara l'insieme dei commi che costituiscono la risposta completa, identificati dal loro \texttt{custom\_id} nel grafo. Le etichette sono usate esclusivamente per la valutazione e non raggiungono mai il modulo di recupero. La Tabella~\ref{tab:composizione_banco} riassume la composizione del banco e il meccanismo che ciascuna categoria mette alla prova.

\begin{table}[H]
    \centering
    \renewcommand{\arraystretch}{1.2}
    \begin{tabular}{|l|c|p{5.2cm}|}
        \hline
        \textbf{Categoria} & \textbf{n} & \textbf{Che cosa mette alla prova} \\
        \hline\hline
        \texttt{SEMANTIC\_DIRECT} & 10 & la domanda è vicina al testo \\
        \hline
        \texttt{SEMANTIC\_INDIRECT} & 10 & la domanda esprime indirettamente il contenuto del testo \\
        \hline
        \texttt{RELATIONAL\_CITES\_LOW\_SIMILARITY} & 10 & rinvii i cui estremi hanno testi poco simili \\
        \hline
        \texttt{RELATIONAL\_CITES\_HIGH\_SIMILARITY} & 10 & rinvii i cui estremi hanno testi molto simili, come condizione di non-effetto atteso \\
        \hline
        \texttt{RELATIONAL\_DEROGATES} & 10 & una regola generale e la disposizione che vi deroga \\
        \hline
        \texttt{MULTI\_ARTICLE} & 10 & commi attesi distribuiti su più articoli \\
        \hline
        \texttt{MULTI\_PARAGRAPH} & 10  & più commi dello stesso articolo \\
        \hline
    \end{tabular}
    \caption{Composizione del benchmark diagnostico.}
    \label{tab:composizione_banco}
\end{table}


\subsubsection{Le due categorie sulla relazione CITES}

Le due categorie dedicate ai rinvii condividono la stessa struttura: il comma sorgente contiene un rinvio esplicito a una disposizione collocata in un altro articolo dello stesso atto. I casi sono suddivisi in due gruppi in base alla similarità semantica tra il comma sorgente e quello richiamato.

La categoria \texttt{RELATIONAL\_CITES\_LOW\_SIMILARITY} comprende i casi in cui i due commi risultano semanticamente distanti. In questa condizione, una ricerca basata esclusivamente sulla similarità semantica può avere difficoltà a recuperarli congiuntamente, poiché il rinvio collega disposizioni che svolgono funzioni differenti o trattano aspetti diversi della stessa disciplina. La presenza di un arco esplicito nel grafo consente invece, una volta individuato il comma sorgente, di raggiungere direttamente la disposizione richiamata, anche quando la somiglianza tra i rispettivi contenuti testuali è ridotta.

Un esempio è rappresentato dal rinvio relativo all'anticipazione delle spese di missione, contenuto nel \emph{Regolamento di contabilità della Provincia di Parma}. L'articolo 74, comma 2, dispone:

\begin{quote}
``L'anticipazione viene corrisposta all'avente diritto secondo le modalità previste al comma 3 dell'art. 73, previa presentazione della copia del relativo atto autorizzativo allo svolgimento della trasferta o missione.''
\end{quote}

La disposizione richiamata, contenuta nell'articolo 73, comma 3, stabilisce invece:

\begin{quote}
``La spesa viene disposta su specifica richiesta firmata dal Dirigente competente tramite l'emissione di appositi buoni economali in cui deve comparire obbligatoriamente l'oggetto della spesa, il suo ammontare, il soggetto percepente la somma, l'impegno di spesa al quale è riferita la spesa.''
\end{quote}


Il primo comma riguarda principalmente l'anticipazione connessa a una missione, mentre il secondo definisce una procedura amministrativo-contabile più generale. La relazione deriva quindi soprattutto dall'importazione, mediante rinvio, di una procedura definita altrove, non da una forte somiglianza tra i due testi.

La categoria \texttt{RELATIONAL\_CITES\_HIGH\_SIMILARITY} comprende invece casi nei quali il rinvio collega disposizioni semanticamente e lessicalmente più vicine. In questa condizione ci si attende un contributo minore dalla struttura a grafo, poiché la disposizione richiamata presenta già caratteristiche che ne favoriscono il recupero attraverso la ricerca vettoriale.

Un esempio riguarda l'utilizzo del mezzo proprio da parte degli amministratori ed è tratto dal \emph{Regolamento per la disciplina dei rimborsi delle spese di viaggio e delle missioni istituzionali degli amministratori della Provincia di Parma}. L'articolo 7, comma 2, afferma:

\begin{quote}
``Nel caso di utilizzo del proprio mezzo, la Provincia di Parma non concorrerà alla copertura assicurativa di detto veicolo e rimarrà comunque esente da ogni responsabilità civile, penale e per infortuni, in relazione ai danni che l'uso del mezzo stesso potrà arrecare a cose o persone, in particolare a terzi, all'Amministratore che utilizza il proprio mezzo, ai trasportati o al mezzo medesimo, come già disciplinato all'art. 5, comma 4, del presente Regolamento.''
\end{quote}

La disposizione richiamata, contenuta nell'articolo 5, comma 4, stabilisce:

\begin{quote}
``Gli Amministratori, autorizzati all'uso del mezzo proprio, devono sollevare l'Amministrazione da qualsiasi responsabilità dichiarando, nell'apposito predisposto modello: [...] che il veicolo utilizzato è in regola con le norme previste per l'assicurazione dei veicoli; [...] di essere in possesso di regolare e valida patente di guida e/o abilitazione alla guida dei veicoli ai sensi del vigente Codice della strada.''
\end{quote}

I due commi condividono il soggetto, la situazione disciplinata e concetti centrali quali l'utilizzo del mezzo proprio, la regolarità del veicolo e l'esclusione della responsabilità dell'Amministrazione. Il rinvio collega pertanto disposizioni caratterizzate da un'elevata vicinanza sia tematica sia lessicale.


\subsubsection{Il trattamento delle altre relazioni giuridiche}

La stessa suddivisione non è stata applicata a \texttt{DEROGATES}. Nei casi presenti nel corpus, questa relazione collega prevalentemente disposizioni che disciplinano lo stesso istituto secondo una struttura regola--eccezione. La disposizione derogante e quella derogata risultano quindi generalmente vicine sul piano tematico e spesso anche su quello lessicale, pur svolgendo funzioni normative differenti.
Un esempio riguarda le auto di servizio disciplinate dal \emph{Regolamento per la gestione del parco veicolare della Provincia di Parma}. La disciplina generale, contenuta nell'articolo 6, commi 7 e 8, prevede:

\begin{quote}
``Le auto di servizio non possono essere utilizzate per il trasporto dall'abitazione all'ufficio o al luogo di lavoro, fatto salvo il personale che svolge servizio di reperibilità.''

``Le stesse dovranno di norma essere stazionate nel luogo di assegnazione corrispondente alla sede dell'Ufficio o del Centro Operativo cui appartiene il dipendente o di quello più vicino alla propria residenza.''
\end{quote}

La disposizione derogante, contenuta nell'articolo 7, comma 1, stabilisce:

\begin{quote}
``Per i dipendenti in servizio di reperibilità è consentito l'utilizzo del mezzo per motivi di servizio in deroga alle disposizioni che prevedono il ricovero presso le sedi della Provincia.''
\end{quote}

Le disposizioni condividono i concetti centrali: le auto di servizio, il personale in reperibilità, l'utilizzo del mezzo e il ricovero presso le sedi della Provincia. La loro affinità semantica non implica tuttavia equivalenza: i commi dell'articolo 6 definiscono la disciplina generale, mentre il comma dell'articolo 7 ne limita l'applicazione introducendo un'eccezione.

Questo esempio evidenzia una differenza rispetto a \texttt{CITES}. Una citazione può collegare sia disposizioni fortemente affini sia disposizioni semanticamente distanti. Nei casi \texttt{DEROGATES} osservati nel corpus, invece, la struttura regola--eccezione tende a mantenere una forte condivisione dell'oggetto normativo. Una suddivisione in condizioni di alta e bassa similarità avrebbe pertanto un interesse sperimentale più limitato e sarebbe ulteriormente indebolita dal numero ridotto di archi \texttt{DEROGATES} disponibili.

Infine, non sono state definite slice analoghe per \texttt{INTERPRETS}, \texttt{REPEALS} e \texttt{REPLACES} poiché, nel grafo utilizzato per la sperimentazione, sono presenti soltanto tre archi \texttt{INTERPRETS}, mentre non risultano archi \texttt{REPEALS} o \texttt{REPLACES}. Tali quantità non consentono di costruire campioni sufficientemente ampi, né di suddividerli ulteriormente in condizioni di alta e bassa similarità. L'introduzione di slice dedicate non avrebbe pertanto prodotto condizioni sperimentali sufficientemente solide per il benchmark.

\subsection{Etichette primarie e di supporto}

Ogni domanda distingue due categorie di commi attesi:

\begin{itemize}
    \item i commi \textbf{primari}, cioè quelli che rispondono direttamente alla domanda e che una ricerca semantica può plausibilmente individuare per vicinanza testuale;
    \item i commi \textbf{di supporto}, cioè le disposizioni necessarie a completare la risposta e la cui pertinenza deriva da un collegamento normativo o strutturale con i commi primari.
\end{itemize}

Questa separazione è lo strumento diagnostico centrale del capitolo, perché consente di osservare separatamente il recupero della risposta diretta e quello del contesto che la completa. La ricerca vettoriale può recuperare entrambe le categorie; l'ipotesi sottoposta a verifica è che l'impiego esplicito delle relazioni produca un guadagno soprattutto sulla copertura dei commi \textit{di supporto}.

\section{Le metriche}
\label{sec:metriche}

Tutte le metriche sono calcolate sul Top-K restituito, ai cutoff 5 e 10. Ciascun cutoff è una selezione a sé: sia i bacini del ramo sia la testa riservata scalano con $k$, quindi un Top-5 non è necessariamente il prefisso di un Top-10

\begin{itemize}
    \item \textbf{Complete HIT@K}: la domanda conta come risolta soltanto se \textit{tutti} i commi attesi si trovano entro il cutoff. È la metrica più severa e la più aderente al caso d'uso, perché una risposta normativa parziale può essere peggiore di nessuna risposta: un utente che riceve la regola generale senza la deroga che la disattiva riceve un'informazione scorretta.
    \item \textbf{Copertura gold@K}: la frazione dei commi attesi presenti nel cutoff. A differenza della precedente premia i progressi parziali.
    \item \textbf{Copertura supporting@K}: la stessa quantità, ristretta ai soli commi di supporto. Il confronto fra S1 e S2 su questa metrica indica quanto l'impiego delle relazioni aiuti a recuperare il contesto che completa la risposta.
    \item \textbf{MRR@K}: il reciproco del rango della prima risposta corretta. Misura la qualità dell'\textit{ordinamento}, non della selezione, ed è la metrica che un sistema generativo a valle consuma effettivamente, perché legge i risultati nell'ordine in cui arrivano.
\end{itemize}

Tutte le medie sono calcolate prima dentro ciascun modello e poi fra i modelli, così che i tre modelli pesino ugualmente. Dove due strategie vengono confrontate, il confronto è \textbf{appaiato}: per ogni valutazione (domanda $\times$ modello) si osserva se il passaggio da S1 a S2 trasformi un risultato incompleto in completo, un risultato completo in incompleto oppure lasci invariato l'esito.

\subsection{Configurazione sperimentale}

Le misure riportate provengono dall'esecuzione \texttt{20260918T180810.030837Z} sul benchmark diagnostico. 

I modelli di embedding confrontati sono i tre già introdotti nella Sezione~\ref{sec:creazione_embedding}, e il capitolo li impiega in \textbf{due ruoli distinti} che vanno dichiarati prima di leggere qualunque tabella che li elenca fianco a fianco, perché altrimenti la stessa riga sembra dire due cose diverse in due punti diversi del capitolo.

\textbf{Primo ruolo, nelle Sezioni~\ref{sec:risultato_principale} e~\ref{sec:costo_grafo}: tre verifiche di robustezza.} La domanda posta lì è se la struttura migliori il recupero, ed è ripetuta con ciascun encoder per verificare che il segno del risultato non dipenda da una sola scelta del modello.

\textbf{Secondo ruolo, nella Sezione~\ref{sec:antitesi}: verifica con encoder contestuale fissato.} Il test mantiene fisso \texttt{pplx-embed-context-v1-0.6b} e misura se il passaggio da S1 a S2 produca ancora un guadagno. Il confronto con \texttt{pplx-embed-v1-0.6b} resta soltanto descrittivo: i due modelli differiscono per addestramento e modalità di costruzione dell'input, quindi non costituiscono un'ablazione a variabile singola del \textit{late chunking}.

\section{Risultati sul benchmark diagnostico}
\label{sec:risultato_principale}

\begin{table}[H]
    \centering
    \renewcommand{\arraystretch}{1.2}
    \begin{tabular}{|l|r|r|r|}
        \hline
        \textbf{Metrica} & \textbf{S1} & \textbf{S2} & \textbf{delta} \\
        \hline
        Complete HIT@10 & 62,4\% & \textbf{79,5\%} & +17,1 pt \\
        Complete HIT@5 & 47,1\% & \textbf{66,2\%} & +19,0 pt \\
        Copertura gold@10 & 80,4\% & \textbf{88,0\%} & +7,6 pt \\
        Copertura gold@5 & 70,0\% & \textbf{79,1\%} & +9,1 pt \\
        Copertura supporting@10 & 69,8\% & \textbf{82,0\%} & +12,2 pt \\
        MRR@10 & \textbf{79,8\%} & 79,5\% & $-$0,4 pt \\
        \hline
    \end{tabular}
    \caption{Confronto complessivo fra ricerca vettoriale pura (S1) e ricerca arricchita dalla struttura (S2), su 70 domande e tre modelli di embedding.}
    \label{tab:risultato_principale}
\end{table}

Sulla composizione scelta per il benchmark, la strategia ibrida ottiene valori superiori su ogni metrica di selezione. Il dato aggregato non va però interpretato separatamente dalle categorie: combina condizioni nelle quali le relazioni possono contribuire e condizioni nelle quali il meccanismo non è richiesto, con pesi stabiliti dalla costruzione del banco. L'unica metrica in cui S2 non migliora è l'MRR, che diminuisce di quattro decimi di punto a causa del costo di ordinamento introdotto dalla fusione.

Sulle 210 valutazioni ottenute combinando 70 domande e tre encoder, il passaggio da S1 a S2 trasforma \textbf{39 risultati incompleti in completi} e \textbf{3 risultati completi in incompleti}; nelle restanti 168 valutazioni l'esito non cambia. Il conteggio rende ispezionabili i casi di disaccordo, ma va letto come statistica descrittiva: le tre valutazioni associate alla stessa domanda sono correlate e non costituiscono osservazioni indipendenti.

\subsection{Il guadagno dipende dalla categoria e non segue interamente l'attesa iniziale}

\begin{table}[H]
    \centering
    \renewcommand{\arraystretch}{1.2}
    \begin{tabular}{|l|r|r|r|}
        \hline
        \textbf{Categoria} & \textbf{S1} & \textbf{S2} & \textbf{delta} \\
        \hline
        \texttt{RELATIONAL\_CITES\_LOW\_SIMILARITY} & 53,3\% & \textbf{90,0\%} & +36,7 pt \\
        \texttt{RELATIONAL\_CITES\_HIGH\_SIMILARITY} & 60,0\% & \textbf{90,0\%} & +30,0 pt \\
        \texttt{RELATIONAL\_DEROGATES} & 70,0\% & \textbf{86,7\%} & +16,7 pt \\
        \texttt{MULTI\_ARTICLE} & 30,0\% & \textbf{70,0\%} & +40,0 pt \\
        \texttt{MULTI\_PARAGRAPH} & 60,0\% & \textbf{66,7\%} & +6,7 pt \\
        \texttt{SEMANTIC\_INDIRECT} & 73,3\% & 73,3\% & 0,0 pt \\
        \texttt{SEMANTIC\_DIRECT} & \textbf{90,0\%} & 80,0\% & $-$10,0 pt \\
        \hline
    \end{tabular}
    \caption{Complete HIT@10 per categoria.}
    \label{tab:per_categoria}
\end{table}

La Tabella~\ref{tab:per_categoria} è più informativa della media che la riassume. Il guadagno non è diffuso: è concentrato sulle categorie costruite attorno a un rinvio o a una distribuzione della risposta su più articoli, cioè dove il meccanismo può agire. Questa regolarità non coincide però interamente con la previsione iniziale: anche \mbox{\texttt{RELATIONAL\_CITES\_HIGH\_SIMILARITY}}, costruita come condizione di non-effetto atteso, guadagna 30 punti di Complete HIT@10. Sulle domande semantiche indirette il grafo è invece \textbf{esattamente neutro}, e sulle domande semantiche dirette \textbf{costa dieci punti}.

La scomposizione del controllo chiarisce il risultato. La copertura dei commi di supporto passa dal 63,3\% al 93,3\%. In nove valutazioni il passaggio da S1 a S2 trasforma un risultato incompleto in completo: S1 aveva già collocato il comma primario nella Top-10, ma non la disposizione richiamata necessaria a completare la risposta; S2 promuove quest'ultima ai ranghi 3--7. Il miglioramento osservato riguarda quindi, in questi casi, il recupero del contesto normativo collegato necessario a completare la risposta. Non si osserva il cambiamento inverso nella categoria.

La previsione di non-effetto non è quindi confermata rispetto alla completezza della risposta. L'alta similarità fra i due estremi dell'arco non garantisce che entrambi compaiano nella Top-10 quando competono con l'intero corpus. Il controllo mette così in evidenza una funzione più circoscritta del grafo: non necessariamente individuare il primo comma pertinente, ma completare il contesto normativo a partire da esso. Questo risultato resta diagnostico, perché la categoria è costruita attorno a rinvii che il sistema può percorrere e non ne misura la frequenza nel traffico reale.

Quest'ultima riga non è un difetto della misura, è la misura funzionante. Le domande dirette sono quelle in cui la ricerca vettoriale già risponde correttamente nel 90\% dei casi; su di esse il grafo non ha nulla da aggiungere, e gli slot che la fusione riserva a candidati di provenienza giuridica sostituiscono risultati vettoriali corretti. È il costo di eseguire un meccanismo dove non serve: la misura complessiva non può essere interpretata come un guadagno netto, ma come un guadagno in condizioni favorevoli e un costo in condizioni sfavorevoli.


\subsection{Il segno del risultato si ripete sui tre encoder}
\label{sec:invarianza_encoder}

\begin{table}[H]
    \centering
    \renewcommand{\arraystretch}{1.2}
    \begin{tabular}{|l|r|r|r|}
        \hline
        \textbf{Modello di embedding} & \textbf{S1} & \textbf{S2} & \textbf{delta} \\
        \hline
        \texttt{embeddinggemma-300m} & 54,3\% & 77,1\% & +22,9 pt \\
        \texttt{pplx-embed-v1-0.6b} & 65,7\% & 82,9\% & +17,1 pt \\
        \texttt{pplx-embed-context-v1-0.6b} & 67,1\% & 78,6\% & +11,4 pt \\
        \hline
    \end{tabular}
    \caption{Complete HIT@10 per modello di embedding, sul banco principale.}
    \label{tab:per_modello}
\end{table}

Ciò che è invariante è il \textbf{segno} del delta: tutti e tre gli encoder migliorano con la struttura, e nessuno peggiora. Ciò che non è invariante è l'\textbf{ampiezza}: il modello che parte dal valore più basso riceve il guadagno maggiore, mentre quello contestuale riceve il guadagno minore (+11,4 punti). La tabella mostra quindi che, nelle condizioni del benchmark, nessuno dei tre encoder rende superfluo il contributo della struttura; non dimostra invece che tale contributo sia indipendente dall'encoder in senso forte.

\section{Il costo della struttura}
\label{sec:costo_grafo}

Un risultato che riporta soltanto i guadagni non è una misura. Nella configurazione valutata il grafo presenta due costi osservabili: uno sull'ordinamento e uno sul tempo di ricerca.

\subsection{Il costo di ordinamento}

Il primo costo è l'MRR: $-$0,4 punti. È piccolo, ma la sua origine è istruttiva. La fusione promuove i commi che compaiono in entrambe le liste, e questi sono tipicamente commi \textit{collegati} alla risposta piuttosto che la risposta stessa; il comma che risponde direttamente può quindi scivolare dal primo posto. La testa riservata esiste per contenere questo effetto, e la sua efficacia è misurabile disattivandola: sul benchmark il rango medio della prima risposta corretta perde oltre cinque punti di MRR, a copertura invariata. È quindi il principale presidio della pipeline contro il deterioramento dell'ordinamento.


\subsection{Il costo computazionale}

Il secondo costo è il tempo. La strategia ibrida aggiunge un percorso di andata e ritorno verso il grafo e la fusione di due graduatorie.

\begin{table}[H]
    \centering
    \renewcommand{\arraystretch}{1.2}
    \begin{tabular}{|l|c|r|r|r|}
        \hline
        \textbf{Strategia} & \textbf{cutoff} & \textbf{mediana} & \textbf{media} & \textbf{p95} \\
        \hline
        S1 & Top-10 & 16 ms & 16 ms & 20 ms \\
        S2 & Top-10 & 18 ms & 21 ms & 35 ms \\
        \hline
    \end{tabular}
    \caption{Latenza di \texttt{strategy.search()}, mediata sui tre modelli. Esclude la codifica della domanda.}
    \label{tab:latenza}
\end{table}

Il sovracosto è di circa 2--5~ms sulla mediana e di 15~ms sul novantacinquesimo percentile. In termini relativi è un incremento sensibile; in condizioni hardware e software della misura, il suo valore assoluto rimane contenuto rispetto alla codifica della domanda e alla successiva generazione della risposta. La valutazione consente quindi di considerare sostenibile il costo computazionale della struttura nella configurazione sperimentata, senza generalizzare la misura ad altri ambienti.

\section{Il contributo del grafo con un encoder contestuale}
\label{sec:antitesi}

Il grafo e il \textit{late chunking} possono entrambi compensare la perdita di contesto prodotta dalla suddivisione di un documento in unità atomiche, ma impiegano segnali differenti. Il \textit{late chunking} incorpora nella rappresentazione di ciascun comma informazioni provenienti dal resto dell'articolo; il grafo codifica invece relazioni esplicite fra disposizioni e le usa durante la ricerca. La parziale sovrapposizione solleva una verifica di robustezza rilevante per il risultato principale: \textit{se la rappresentazione vettoriale è già contestualizzata, le relazioni esplicite conservano un contributo osservabile?}

\begin{table}[H]
    \centering
    \renewcommand{\arraystretch}{1.2}
    \begin{tabular}{|l|r|r|r|}
        \hline
        \textbf{Metrica con encoder contestuale} & \textbf{S1} & \textbf{S2} & \textbf{delta} \\
        \hline
        Complete HIT@10 & 67,1\% & \textbf{78,6\%} & +11,4 pt \\
        Copertura gold@10 & 81,3\% & \textbf{86,8\%} & +5,5 pt \\
        Copertura supporting@10 & 71,9\% & \textbf{79,7\%} & +7,8 pt \\
        MRR@10 & \textbf{75,2\%} & 74,8\% & $-$0,5 pt \\
        \hline
    \end{tabular}
    \caption{Contributo della struttura mantenendo fisso l'encoder contestuale.}
    \label{tab:antitesi_test_diretto}
\end{table}

Sull'intero banco di 70 domande, la struttura guadagna 11,4 punti di Complete HIT@10 anche quando i commi sono rappresentati mediante l'encoder contestuale. La copertura dei commi gold aumenta di 5,5 punti e, sulle domande per le quali sono definiti commi di supporto, la loro copertura cresce di 7,8 punti. L'MRR diminuisce invece di 0,5 punti, riproducendo su scala ridotta il costo di ordinamento osservato nel confronto aggregato. Il risultato respinge quindi, nelle condizioni valutate, l'ipotesi forte secondo cui la contestualizzazione dell'encoder assorbirebbe interamente il contributo delle relazioni esplicite.

Il risultato aggregato non chiarisce tuttavia in quali condizioni il grafo mantenga il proprio vantaggio. La scomposizione per categoria mostra che l'effetto non consiste in un miglioramento uniforme.

\begin{table}[H]
    \centering
    \renewcommand{\arraystretch}{1.2}
    \begin{tabular}{|l|r|r|r|}
        \hline
        \textbf{Categoria} & \textbf{S1} & \textbf{S2} & \textbf{delta} \\
        \hline
        \texttt{RELATIONAL\_CITES\_LOW\_SIMILARITY} & 50,0\% & \textbf{80,0\%} & +30,0 pt \\
        \texttt{RELATIONAL\_CITES\_HIGH\_SIMILARITY} & 50,0\% & \textbf{80,0\%} & +30,0 pt \\
        \texttt{RELATIONAL\_DEROGATES} & 90,0\% & \textbf{100,0\%} & +10,0 pt \\
        \texttt{MULTI\_ARTICLE} & 50,0\% & \textbf{70,0\%} & +20,0 pt \\
        \rowcolor{orange!25}
        \texttt{MULTI\_PARAGRAPH} & 70,0\% & 70,0\% & 0,0 pt \\
        \texttt{SEMANTIC\_INDIRECT} & 80,0\% & 80,0\% & 0,0 pt \\
        \texttt{SEMANTIC\_DIRECT} & \textbf{80,0\%} & 70,0\% & $-$10,0 pt \\
        \hline
    \end{tabular}
    \caption{Complete HIT@10 per categoria mantenendo fisso l'encoder contestuale.}
    \label{tab:antitesi_per_categoria}
\end{table}

Il vantaggio di S2 si concentra nelle categorie relazionali e in \texttt{MULTI\_ARTICLE}; \texttt{MULTI\_PARAGRAPH} e \texttt{SEMANTIC\_INDIRECT} rimangono invarianti, mentre \texttt{SEMANTIC\_DIRECT} perde dieci punti. Questa distribuzione è coerente con quanto osservato nel confronto per categoria (Tabella \ref{tab:per_categoria}), tranne che per \texttt{MULTI\_PARAGRAPH}: nel confronto precedente beneficia della struttura, ma con l'encoder contestuale non mostra alcun guadagno aggiuntivo.

Fra le condizioni neutrali, \texttt{MULTI\_PARAGRAPH} costituisce il controllo più direttamente pertinente al rapporto fra grafo e \textit{late chunking}: richiede di recuperare più commi dello stesso articolo, ossia informazioni comprese nella medesima finestra di contestualizzazione dell'encoder. Se il grafo compensasse principalmente la perdita di contesto intra-articolo, il suo contributo dovrebbe quindi ridursi, o scomparire, in questa categoria.

\begin{table}[H]
    \centering
    \renewcommand{\arraystretch}{1.2}
    \begin{tabular}{|l|r|r|r|}
        \hline
        \textbf{Metrica su \texttt{MULTI\_PARAGRAPH}} & \textbf{S1} & \textbf{S2} & \textbf{delta} \\
        \hline
        Complete HIT@10 & 70,0\% & 70,0\% & 0,0 pt \\
        Copertura gold@10 & 76,7\% & 76,7\% & 0,0 pt \\
        Copertura supporting@10 & 75,0\% & 75,0\% & 0,0 pt \\
        MRR@10 & 65,8\% & 65,8\% & 0,0 pt \\
        \hline
    \end{tabular}
    \caption{Confronto fra S1 e S2 sulle dieci domande \texttt{MULTI\_PARAGRAPH}, mantenendo fisso l'encoder contestuale.}
    \label{tab:multi_paragraph_context}
\end{table}

S1 e S2 ottengono gli stessi risultati su tutte le metriche considerate. In questa categoria non emerge quindi un vantaggio ulteriore del grafo. Il pareggio è coerente con una sovrapposizione dei due segnali quando il contesto necessario è interno allo stesso articolo; la Tabella~\ref{tab:antitesi_per_categoria} mostra però che il grafo conserva un contributo nelle altre categorie, ossia quelle che coinvolgono relazioni esplicite o più articoli.

Le due tecniche operano su forme di contesto differenti. Il \textit{late chunking} contestualizza ciascun comma mediante le informazioni contenute nella finestra ricevuta dall'encoder, che in questa implementazione coincide con l'articolo. La ricerca vettoriale continua a interrogare l'intero corpus, ma ogni rappresentazione incorpora direttamente soltanto il contesto intra-articolo.

Il grafo codifica invece relazioni esplicite tra disposizioni, indipendentemente dal fatto che appartengano o meno allo stesso articolo. Durante la ricerca, tali relazioni possono essere impiegate per promuovere candidati presenti nel bacino vettoriale. Il segnale strutturale può quindi estendersi oltre il perimetro contestualizzato dall'encoder.

I risultati sono coerenti con questa distinzione. In \texttt{MULTI\_PARAGRAPH}, dove i commi attesi appartengono allo stesso articolo, l'aggiunta del grafo non produce un beneficio osservabile. Il vantaggio rimane invece nelle categorie relazionali e in \texttt{MULTI\_ARTICLE}, cioè in condizioni nelle quali la risposta può coinvolgere collegamenti o contenuti distribuiti oltre il singolo articolo. Nel benchmark considerato emerge quindi una parziale sovrapposizione sul contesto intra-articolo, ma non una completa ridondanza tra contestualizzazione e relazioni esplicite.

\section{Limiti della valutazione}

Prima di sintetizzare i risultati complessivi, è opportuno chiarire ciò che le misure adottate non consentono di stabilire.

\textbf{La dimensione del campione è modesta.} Settanta domande consentono di descrivere differenze ampie, ma non di attribuire un significato forte a scarti di pochi punti. Le 210 combinazioni domanda $\times$ modello non sono osservazioni indipendenti, perché i tre encoder vengono valutati sugli stessi 70 quesiti.

\textbf{Il benchmark è costruito, non campionato.} Cinque categorie su sette sono state scritte attorno a un meccanismo del sistema. I risultati mostrano quindi che cosa accade nelle condizioni rappresentate dalle categorie, ma non stimano la frequenza di tali condizioni nel traffico reale.

\textbf{Il corpus è uno solo.} Le proporzioni fra i tipi di relazione di questo grafo --- 1.030 archi \texttt{CITES} contro 41 \texttt{DEROGATES} e 3 \texttt{INTERPRETS}, con \texttt{REPEALS} e \texttt{REPLACES} assenti --- determinano quanto il ramo giuridico possa contribuire. Un corpus con una fitta trama di abrogazioni e sostituzioni si comporterebbe diversamente, e presumibilmente in favore della struttura, ma questo non è stato misurato.

\textbf{La qualità della risposta generata non è misurata.} Tutto ciò che questo capitolo misura è il recupero: quali commi arrivano, in quale ordine. Che un sistema generativo a valle produca con quei commi una risposta corretta è una proprietà distinta, che richiederebbe un impianto di valutazione diverso.

\section{Sintesi del capitolo}

Considerando tutte le 70 domande, la ricerca arricchita dal grafo porta la Complete HIT@10 dal 62,4\% al 79,5\%, con un guadagno di 17,1 punti. Sulle 210 valutazioni ottenute combinando le domande con i tre encoder, il passaggio da S1 a S2 trasforma 39 risultati incompleti in completi e 3 risultati completi in incompleti; nelle restanti 168 valutazioni l'esito non cambia. Il vantaggio riguarda la completezza della selezione, mentre l'MRR diminuisce di 0,4 punti: il grafo recupera più spesso l'intera risposta, ma può modificare l'ordine dei primi risultati.

La verifica con encoder contestuale fissato conduce alla stessa conclusione generale: S2 supera S1 di 11,4 punti di Complete HIT@10, di 5,5 punti nella copertura dei commi gold e di 7,8 punti nella copertura dei commi di supporto, a fronte di una diminuzione di 0,5 punti di MRR. Il vantaggio si concentra nelle categorie relazionali e in \texttt{MULTI\_ARTICLE}; \texttt{MULTI\_PARAGRAPH} e \texttt{SEMANTIC\_INDIRECT} rimangono invarianti, mentre \texttt{SEMANTIC\_DIRECT} peggiora. Il \textit{late chunking} e il grafo risultano quindi parzialmente sovrapposti sul contesto intra-articolo, ma non ridondanti rispetto all'insieme delle relazioni rappresentate nel benchmark.

Nel complesso, l'esperimento mostra che la struttura può aumentare la completezza del recupero nelle condizioni per cui le relazioni risultano pertinenti, con un costo contenuto ma osservabile sull'ordinamento. Le conclusioni descrivono il comportamento dei meccanismi nelle condizioni rappresentate dal benchmark e non costituiscono una stima delle prestazioni attese su domande reali.